\documentclass{beamer}
\usepackage{mathtools}

\newtheorem{algorithm}{Algorithm}
\newtheorem{setup}{Setup}
\newtheorem{encryption}{Encryption}
\newtheorem{decryption}{Decryption}
\newcommand{\zmod}[1]{\ensuremath{\mathbb{Z}/#1\mathbb{Z}}}

\usetheme{PaloAlto}
\usecolortheme{penguin}

\title[Elliptic Curve Cryptography]{Elliptic Curves and
    Their Applications in Cryptography}
\author{Tim Shaffer\thanks{Adviser: Dr. Jacek Fabrykowski}}
\institute[YSU]{Youngstown State University}
\date{MathFest 2014}

\begin{document}
    \frenchspacing
    \frame{\maketitle}

    \begin{frame}
        \frametitle{Fermat's Last Theorem}
        Pierre de Fermat: % remember title of book (Arithmetica)
        \begin{quote}
            It is impossible to write a cube as the sum of two cubes, a fourth
            power as the sum of two fourth powers, and, in general, any power
            beyond the second as the sum of two similar powers. For this I
            have discovered a truly wonderful proof but the margin is too
            small to contain it.
        \end{quote}
        \vfill
        1995--Andrew Wiles published the first successful proof.
    \end{frame}

    \section{Definition}
    \begin{frame}
        \begin{definition}
            An elliptic curve is a projective algebraic curve with
            affine coordinates given by \[y^2 = x^3 + \alpha x + \beta\]
            where the discriminant
            \[\Delta = -16(4\alpha^3 + 27\beta^2) \neq 0.\]
        \end{definition}
        \hyperlink{definitions}{\beamergotobutton{Definitions}}
        \vfill
        Elliptic curves over finite fields (usually \zmod{p}) are of
        particular interest in cryptography.
    \end{frame}

    \begin{frame}
        \frametitle{Elliptic Curves}
        \centering
        \includegraphics[height=.9\textheight]{EllipticCurveCatalog.pdf}
    \end{frame}

    \begin{frame}
        \frametitle{Elliptic Curve \(y^2 = x^3 - x\) on \zmod{61}}
        \centering
        \includegraphics[width=\textwidth]{Z61.pdf}
    \end{frame}

    \section{Group Properties}
    \subsection{Additive Operation}
    \begin{frame}
        \frametitle{Additive Operation}
        \hyperlink{points}{\beamergotobutton{Point at Infinity}}
        \centering
        \includegraphics[width=\textwidth]{ECClines.pdf}
    \end{frame}

    \begin{frame}
        \frametitle{Additive Operation}
        \begin{theorem}
            For an elliptic curve \(\mathcal{E}(\mathbb{F}_q)\)
            and \(a, b, c, d \in \mathbb{F}_q\),
            let \(P = (a, b)\) and \(Q = (c, d)\) be two points
            on \(\mathcal{E}\) such that \(Q \neq -P = (a, -b)\). Define
            \[m =
                \begin{dcases}
                    \frac{d - b}{c - a} & \textnormal{if}\ a \neq c \\
                    \frac{3a^2 + \alpha}{2b} & \textnormal{if}\ a = c.
                \end{dcases}\]
            Then the point \(P + Q\) is given by \(R = (g, h)\) where
            \begin{align*}
                g &= m^2 - a - c \\
                h &= ma - mg - b.
            \end{align*}
        \end{theorem}
        \hyperlink{add}{\beamergotobutton{Proof}}
    \end{frame}

    \subsection{ECDLP}
    \begin{frame}
        \frametitle{Elliptic Curve Discrete Logarithm Problem}
        Given \(P = (a, b)\) on \(\mathcal{E}\), it is possible to
        efficiently compute \(\overbrace{P + \dots + P}^{n\ \textrm{times}}\),
        denoted \([n]P\).

        \hyperlink{mul}{\beamergotobutton{Implementation}}
        \vfill
        However, given \(P\) and \([n]P\), it can be \emph{very} difficult to
        compute the value of \(n\).
    \end{frame}

    \section{Cryptography}
    \begin{frame}
        Assume that all parties agree in advance on a choice of elliptic curve
        \(\mathcal{E}\) on finite field \(\mathbb{F}_p\),
        \(P \in \mathcal{E}(\mathbb{F}_p)\), and paring \(\omega\)
        on \(\mathcal{E}\). \hyperlink{pair}{\beamergotobutton{Pairings}}
        \vfill
        The chosen parameters are assumed to be public knowledge and to
        possess the properties appropriate for security.
        \vfill
        Let \(M \in \mathcal{E}(\mathbb{F}_p)\) be a message (encoded as
        a point on \(\mathcal{E}(\mathbb{F}_p)\))
        that Alice would like to send to Bob.
    \end{frame}

    \subsection{ECDH}
    \begin{frame}
        \frametitle{Elliptic Curve Diffie-Hellman Exchange}
        \begin{setup}
            \begin{itemize}
                \item Alice chooses a random secret \(a \in \mathbb{Z}\) and
                    sends \(A = [a]P\) to Bob over an insecure channel.
                \item In the same way, Bob chooses a random secret
                \(b \in \mathbb{Z}\) and sends \(B = [b]P\) to Alice.
            \end{itemize}
        \end{setup}
        \vfill
        \begin{algorithm}
            \begin{itemize}
                \item Alice computes \(Q = [a]B = [ab]P\).
                \item Bob computes \(Q = [b]A = [ba]P\).
            \end{itemize}
        \end{algorithm}
    \end{frame}

    \subsection{Elgamal}
    \begin{frame}
        \frametitle{Elgamal Encryption}
        \begin{setup}
            Bob chooses a private key \(b \in \mathbb{Z}\) and computes
            his public key \(B = [b]P\). Bob is free to publish \(B\).
        \end{setup}
        \vfill
        \begin{encryption}
            Alice chooses a random \(k \in \mathbb{Z}\)
            and computes \(C_1 = [k]P\),
            \(C_2 = M + [k]B\) and sends \((C_1, C_2)\) to Bob.
        \end{encryption}
        \vfill
        \begin{decryption}
            Bob computes \(C_2 - [b]C_1 = M + [k]B - [b][k]P =
            M + [kb]P - [bk]P = M\).
        \end{decryption}
    \end{frame}

    \begin{frame}
        Assume that all parties choose in advance the
        following parameters.
        \begin{itemize}
            \item \(\mathcal{E}\): \(y^2 = x^3 - x + 1\)
            \item \(\mathbb{F_p}\): \zmod{113}
            \item \(P\): \((69, 96)\)
        \end{itemize}
        \vfill
        Let \(M = (53, 111)\) be a block of a secret message.
    \end{frame}

    \begin{frame}
        \frametitle{Elgamal Example}
        \begin{setup}
            Bob chooses a private key \(b = 8\) and publishes
            his public key \(B = [b]P = [8](69, 96) = (95, 17)\).
        \end{setup}
        \vfill
        \begin{encryption}
            Alice chooses a random \(k = 11\)
            and computes \(C_1 = [k]P = [11](69, 96) = (71, 99)\),
            \(C_2 = M + [k]B = (53, 111) + [11](95, 17) = (96, 23)\)
            and sends \((C_1, C_2)\) to Bob.
        \end{encryption}
        \vfill
        \begin{decryption}
            Bob computes \(C_2 - [b]C_1 = (96, 23) - [8](71, 99) = (53, 111)\).
        \end{decryption}
    \end{frame}

%   \subsection{Tripartite Key Exchange}
%   \begin{frame}
%       \frametitle{Tripartite Key Exchange}
%       \begin{setup}
%           Alice, Bob, and Carl each choose a random secret
%           \(a, b, c \in \mathbb{Z}\), respectively. They compute and
%           share \(A = [a]P\), \(B = [b]P\), and \(C = [c]P\)
%           over insecure channels.
%       \end{setup}
%       \vfill
%       \begin{algorithm}
%           Each computes the shared secret as follows using his/her
%           respective secret
%           \begin{itemize}
%               \item Alice: \(\omega(B, C)^a\)
%               \item Bob: \(\omega(A, C)^b\)
%               \item Carl: \(\omega(A, B)^c\)
%           \end{itemize}
%           since \(\omega(B, C)^a = \omega(A, C)^b = \omega(A, B)^c =
%               \omega(P, P)^{abc}\)
%       \end{algorithm}
%   \end{frame}

    \subsection{ID-based Encryption}
    \begin{frame}
        \frametitle{Identity-based Encryption}
        Let Trent act as the trusted authority.
        \vfill
        \begin{setup}
            \begin{enumerate}
                \item Trent chooses a master secret \(s \in \mathbb{Z}\) and
                    publishes \(S = [s]P\).
                \item Bob encodes his identity (e.g. username, email, etc.)
                    as \(b \in \mathbb{Z}\). Anyone can compute Bob's public
                    key \(B = [b]P\).
                \item Bob requests his private key
                    \(E = [s]B = [sb]P\) from Trent.
            \end{enumerate}
        \end{setup}
    \end{frame}

    \begin{frame}
        \frametitle{Identity-based Encryption}
        \begin{encryption}
            Alice chooses a random secret \(t \in \mathbb{Z}\) and sends
            \((U, V) = ([t]P, M + \omega(B, S)^t)\) to Bob.
        \end{encryption}
        \vfill
        \begin{decryption}
            Observe that \[\omega(E, U) = \omega([s]B, [t]P) =
            \omega(B, P)^{st} = \omega(B, [s]P)^t = \omega(B, S)^t.\]
            Bob computes \(M = V - \omega(E, U) = V - \omega(B, S)^t\).
        \end{decryption}
    \end{frame}



    \section{Comparisons}
    \begin{frame}
        \frametitle{Security}
        \centering
        NIST Recommended Key Sizes\\
        \vspace{.5cm}
        \begin{tabular}{c c c}
            Symmetric Key Size & RSA and DH Key Size & EC Key Size \\ \hline
            80 & 1024 & 160 \\
            112 & 2048 & 224 \\
            128 & 3072 & 256 \\
            192 & 7680 & 384 \\
            256 & 15360 & 521 \\
        \end{tabular}
    \end{frame}

    \begin{frame}
        \frametitle{Efficiency}
        \centering
        Operations per 10 seconds \\
        \vspace{.5cm}
        \begin{tabular}{c|c c}
             & Sign & Verify \\ \hline \hline
            RSA 1024 & 51670 & 819461 \\
            RSA 2048 & 7669 & 247067 \\
            RSA 4096 & 1072 & 66025 \\
            \hline
            ECDSA 160 & 163316 & 44279 \\
            ECDSA 224 & 144934 & 64201 \\
            ECDSA 521 & 31128 & 13824 \\
        \end{tabular}
    \end{frame}

    \section{References}
    \begin{frame}[allowframebreaks]
        \frametitle{References}
        \begin{itemize}
            \item Baker, Alan. \emph{A Comprehensive Course in Number Theory}.
                New York: Cambridge UP, 2012. Print.
            \item Bernstein, Daniel J. and Tanja Lange. \emph{SafeCurves:
                choosing safe curves for elliptic-curve cryptography}.
                http://safecurves.cr.yp.to, accessed 4 August 2014.
            \item Boutet, Emmanuel. Example Elliptic Curves. Digital image.
                \emph{Wikimedia Commons}.
                Wikimedia Foundation, 25 Oct. 2007.
                Web. 16 July 2014. (CC BY-SA 3.0)
            \item Ireland, Kenneth F., and Michael I. Rosen. \emph{A Classical
                Introduction to Modern Number Theory}. Vol. 84. New York:
                Springer-Verlag, 1990. Print. Graduate Texts in Mathematics.
            \item Ling, San, Huaxiong Wang, and Chaoping Xing. \emph{Algebraic
                Curves in Cryptography}. Boca Raton: CRC, 2013. Print.
            \item Menezes, Alfred J., Paul C. Van Oorschot, and
                Scott A. Vanstone. \emph{Handbook of Applied Cryptography}.
                Boca Raton: CRC, 1997. Print.
            \item ``The Case for Elliptic Curve Cryptography.''
                \emph{Central Security Service}.
                NSA, 15 Jan. 2009. Web. 06 Aug. 2014.
        \end{itemize}
    \end{frame}

    \appendix
    \section{Definitions}
    \begin{frame}[label=definitions]
        \frametitle{Affine Space}
        \begin{definition}
            Given a finite field with \(q\) elements, \(\mathbb{F}_q\),
            affine \(n\)-space over \(\mathbb{F}_q\), denoted
            \(A^n(\mathbb{F}_q)\), is the set of
            \(n\)-tuples \((a_1, a_2, \dots, a_n)\) with
            \(a_i \in \mathbb{F}_q.\)
        \end{definition}
        \vfill
        \begin{definition}
            A \textbf{point} in \(A^n(\mathbb{F}_q)\) is an \(n\)-tuple
            \((a_1, a_2, \dots, a_n)\) for \(a_i \in \mathbb{F}_q\).
        \end{definition}
    \end{frame}

    \begin{frame}
        \frametitle{Projective Space}
        \begin{definition}
            Projective \(n\)-space over \(\mathbb{F}_q\), denoted
            \(P^n(\mathbb{F}_q)\), is the set of equivalence classes of
            nonzero elements of \(A^{n+1}(\mathbb{F}_q)\) under the
            equivalence relation
            \[(a_0, a_1, \dots, a_n) \sim (b_0, b_1, \dots, b_n)\] iff there
            exists a \(0 \neq \lambda \in \mathbb{F}_q\) such that
            \[a_i = \lambda b_i\] for all \(i = 0, 1, \dots, n\).
        \end{definition}
    \end{frame}

    \begin{frame}
        \frametitle{Projective Space}
        \begin{definition}
            A \textbf{point} in \(P^n(\mathbb{F}_q)\), denoted
            \([a_0, a_1, \dots, a_n]\), is the equivalence class
            containing \((a_0, a_1, \dots, a_n)\).
        \end{definition}
        \vfill
        While \(A^2(\mathbb{F}_q)\) has \(q^2\) points,
        \(P^2(\mathbb{F}_q)\) has \(q^2 + q + 1\) points.
        \vfill
        The points in \(P^2(\mathbb{F}_q)\) can be broken into 2 subsets:
        \begin{itemize}
            \item \(q^2\) \textbf{finite points} of the form \([1, a_1, a_2]\)
            that map to \(A^2(\mathbb{F}_q)\)
            \item \(q+1\) \textbf{points at infinity} of the
            form \([0, a_0, a_1]\)
            with the structure of \(P^1(\mathbb{F}_q)\)
        \end{itemize}
    \end{frame}

    \begin{frame}
        \frametitle{Algebraic Curves}
        \begin{definition}
            An affine algebraic curve over \(\mathbb{F}_q\) is defined by
            \(f(x, y) = 0\) for an irreducible polynomial
            \(f(x, y) \in \mathbb{F}_q[x, y]\).
        \end{definition}
        \vfill
        \begin{definition}
            A projective algebraic curve over \(\mathbb{F}_q\) is defined by
            \(f(x, y, z) = 0\) for an irreducible homogeneous polynomial
            \(f(x, y, z) \in \mathbb{F}_q[x, y, z]\).
        \end{definition}
    \end{frame}

    \begin{frame}
        \frametitle{Singular Points}
        \begin{definition}
            A point \(P\) on an affine curve \(f(x, y) = 0\) is called
            \textbf{singular} if
            \[\left(\frac{\partial f}{\partial x}(P),
            \frac{\partial f}{\partial y}(P)\right) = (0, 0)\]
        \end{definition}
        \vfill
        \begin{definition}
            An algebraic curve is called \textbf{nonsingular} or
            \textbf{smooth} if it contains no singular points.
        \end{definition}
    \end{frame}

    \begin{frame}
        \frametitle{Elliptic Curves}
        The \textbf{projective Weierstrass equation} of an elliptic curve is
        given by
        \[y^2z + \alpha_1xyz + \alpha_3yz^2 =
        x^3 + \alpha_2x^2z + \alpha_4xz^2 + \alpha_6z^3.\]
        \vfill
        The affine \textbf{Weierstrass normal form} of an elliptic curve is
        \[y^2 - x^3 - \alpha x - \beta = 0\] with
        discriminant \(\Delta = -16(4\alpha^3 + 27\beta^2)\).
    \end{frame}

    \section{Points in \(P^2(\mathbb{F}_q)\)}
    \begin{frame}[label=points]
        \frametitle{Number of Points in \(P^2(\mathbb{F}_q)\)}
        It is clear that the number of points in \(A^n(\mathbb{F}_q)\)
        is \(q^n\).
        \vfill
        \begin{proof}
            Observe that
            \[P^2(\mathbb{F}_q) = \{[1, a, b] | a, b \in \mathbb{F}_q\}
                \cup \{[0, 1, a] | a \in \mathbb{F}_q\}
                \cup \{[0, 0, 1]\}\]
            It follows that the number of points in \(P^2(\mathbb{F}_q)\)
            is \[q^2 + q + 1.\]
        \end{proof}
    \end{frame}

    \section{Maps Between Spaces}
    \begin{frame}
        \frametitle{Mappings Between Projective and Affine Spaces}
        Let \(K\) be a field and define \(H\) as the points at infinity
        of \(P^n(K)\). Mappings \(\lambda\) and \(\phi\) between
        affine and projective spaces are defined as follows.
        \begin{align*}
            \lambda&: A^n(K) \to P^n(K) \\
            &\lambda(a_1, a_2, \dots, a_n) = [1, a_1, a_2, \dots, a_n]
        \end{align*}
        \begin{align*}
            \phi&: P^n(K) - H \to A^n(K) \\
            &\phi([b_0, b_1, \dots, b_n]) =
                \left(\frac{b_1}{b_0}, \frac{b_2}{b_0}, \dots,
                \frac{b_n}{b_0}\right)
        \end{align*}
    \end{frame}

    \section{Addition on \(\mathcal{E}\)}
    \begin{frame}[label=add]
        \frametitle{Correctness of Additive Operation on \(\mathcal{E}\)}
        Let \(\ell\), the line passing through \(P\) and \(Q\), be given by
        \(y = mx - C\) with \(m\) as defined for the addition
        operation and \(k = b - ma\). Also let \(S = (g', h')\) be the third
        intersection of
        \(\mathcal{E}\) and \(\ell\). Substituting the equation for \(\ell\)
        into that of \(\mathcal{E}\),
        \[(mx + k)^2 = x^3 + \alpha x + \beta\] which expands to
        \[f(x) = x^3 - m^2x^2 + (\alpha - 2mk)x + \beta - k^2 = 0.\]
    \end{frame}
    \begin{frame}
        \frametitle{Correctness of Additive Operation on \(\mathcal{E}\)}
        Since \(a, c, g'\) are the \(x\) coordinates of \(p, Q, S\),
        respectively,
        \[f(x) = (x - a)(x - c)(x - g')\]
        and by expanding and comparing coefficients,
        \begin{align*}
            g' &= M^2 - a - c \\
            h' &= Mg' + N
        \end{align*}
        so
        \[P + Q = S = (g', -h').\]
    \end{frame}

    \section{Pairings}
    \begin{frame}[label=pair]
        \frametitle{Pairings on Elliptic Curves}
        \begin{definition}
            Given abelian groups \(G_1, G_2, G_3\), a pairing
            \(\omega: G_1 \times G_2 \to G_3\) maps every pair of elements
            in \(G_1 \times G_2\) to some element in \(G_3\).
        \end{definition}
        \vfill
        A \emph{cryptographically useful} pairing is also
        \begin{itemize}
            \item bilinear: if \(g_1, g_1' \in G_1\)
            and \(g_2, g_2' \in G_2\)
            then \(\omega(g_1g_1', g_2) =
            \omega(g_1, g_2)\omega(g_1', g_2)\)
            and \(\omega(g_1, g_2g_2') =
            \omega(g_1, g_2)\omega(g_1, g_2')\).
            \item nondegenerate: if \(\omega(g_1, g_2) = 1\)
            for all \(g_2 \in G_2\) then it follows that \(g_1 = 1\).
        \end{itemize}
    \end{frame}

    \section{Implementation}
    \begin{frame}[label=mul]
        \frametitle{Efficient Computation of \([n]P\)}
        \begin{algorithm}
            \begin{enumerate}
                \item Write \(n\) in its binary form, i.e.
                \(n = n_0 + 2n_1 + 2^2n_2 + \dots + 2^tn_t\),
                with \(n_i \in \{0, 1\}\) and \(n_t = 1\).
                \item Let \(P_0 = P\).
                \item For \(i = 1, \dots, t\), compute \(P_i = [2^i]P =
                [2]P_{i - 1}\) recursively.
                \item Then, \([m]P =\sum _{i = 0} ^t [n_i]P_i\).
            \end{enumerate}
        \end{algorithm}
        \vfill
        While na\"{i}ve application of the group operator
        requires \(n\) additions, this algorithm can be carried out
        in \(2t \leq 2 \log n\) additions.
    \end{frame} % brush up on \rho, baby-step giant-step

    \subsection{Order of \(\mathcal{E}(\mathbb{F}_q)\)}
    \begin{frame}
        \frametitle{Order of \(\mathcal{E}(\mathbb{F}_q)\)}
        A curve with a small number of points would be vulnerable to
        cryptanalysis.
        \vfill
        If an elliptic curve on \(\mathbb{F}_p\) has exactly \(p\) points,
        the ECDLP can be transformed into addition in \(\mathbb{Z}_p\).
        \vfill
        \begin{definition}[Hasse-Weil Bound]
            Let \(N\) be the number of points in \(\mathbb{F}_q\)
            on an elliptic curve \(\mathcal{E}\).
            \[q + 1 - 2\sqrt{q} \leq N \leq q + 1 + 2\sqrt{q}.\]
        \end{definition}
    \end{frame}
    
    \subsection{Schoof's Algorithm}
    \begin{frame}
        \frametitle{Schoof's Algorithm}
        The number of points on an elliptic curve can be calculated
        efficiently using a deterministic polynomial time algorithm.
        \vfill
        Schoof's Algorithm works by computing \(q + 1 - N \pmod p\) for
        a large number of primes whose product is greater than \(4\sqrt{q}\),
        then calculating \(q + 1 - N\) by the Chinese Remainder Theorem.
        \vfill
        With improvements by Atkin and Elkies, Schoof's Algorithm
        runs in \(O(\log^4 q)\) time.
    \end{frame}

    \section{Sample Curve}
    \begin{frame}
        \frametitle{secp224r1}
        \[y^2 = x^3 -3x + \beta\]
        over \zmod{p} with
        \begin{align*}
        p &= 2^224 - 2^96 + 1 \\
        \beta &= b4~05~0a~85~0c~04~b3~ab~f5~41~32~56~50 \\
        & 44~b0~b7~d7~bf~d8~ba~27~0b~39~43~23~55~ff~b4 \\
        \end{align*}
        and base point
        \begin{align*}
        (b7~0e~0c~bd~6b~b4~bf~7f~32~13~90~b9~4a~03 \\
        c1~d3~56~c2~11~22~34~32~80~d6~11~5c~1d~21, \\
        bd~37~63~88~b5~f7~23~fb~4c~22~df~e6~cd~43 \\
        75~a0~5a~07~47~64~44~d5~81~99~85~00~7e~34) \\
        \end{align*}
    \end{frame}

    \section{ECDH Example}
    \begin{frame}
        \frametitle{ECDH Example}
        \begin{setup}
            \begin{itemize}
                \item Alice chooses \(a = 7\) and sends \(A = [a]P =
                    [7](69, 96) = (62, 96)\) to Bob.
                \item Bob chooses \(b = 12\) and sends \(B = [b]P =
                    [12](69, 96) = (60, 87)\) to Alice.
            \end{itemize}
        \end{setup}
        \vfill
        \begin{algorithm}
            \begin{itemize}
                \item Alice computes \([a]B = [7](60, 87) = (67, 2)\).
                \item Bob computes \([b]A = [12](62, 96) = (67, 2)\).
            \end{itemize}
        \end{algorithm}
    \end{frame}

    \section{Factorization}
    \begin{frame}
        \frametitle{Elliptic Curve Factorization}
        Pollard's~\(p - 1\)~algorithm can find prime factors \(p\) of a
        composite integer for which \(p - 1\) is smooth with respect to some
        relatively small bound \(k\).
        \vfill
        \begin{definition}
            An integer is called \textbf{\(k\)-smooth} if all of its prime
            factors are less than \(k\).
            \end{definition}
            \vfill
            Elliptic curve factorization is a generalization of
            Pollard's~\(p - 1\)~algorithm using random elliptic curve groups
            over \zmod{p}.
            \end{frame}

    \begin{frame}
        \centering
        \begin{columns}[T]
            \begin{column}[T]{.2\textwidth}
                \centering
                \includegraphics[width=\textwidth]{cc.pdf}
            \end{column}
            \begin{column}[T]{.2\textwidth}
                \centering
                \includegraphics[width=\textwidth]{by.pdf}
            \end{column}
            \begin{column}[T]{.2\textwidth}
                \centering
                \includegraphics[width=\textwidth]{sa.pdf}
            \end{column}
        \end{columns}
        \vfill
        This work is licensed under the
        Creative~Commons Attribution-ShareAlike~4.0 International~License.
        To view a copy of this license, visit
        \url{http://creativecommons.org/licenses/by-sa/4.0/}.
    \end{frame}
\end{document}
